\documentclass[10pt,a4paper]{amsart}
\usepackage[margin=19mm]{geometry}
\usepackage{amsmath,amssymb,mathtools}
\usepackage[scr=rsfs,cal=euler]{mathalfa}
\usepackage{xfrac}
\usepackage[unicode,hidelinks,bookmarks=false]{hyperref}
\newcommand{\ie}{i.e.\ }
\newcommand{\cf}{cf.\ }
\newcommand{\resp}{respectively\ }
\newcommand{\Subsection}{\subsection}
\providecommand{\nobreakdash}{\nobreak\hbox{-}}
\setlength{\emergencystretch}{2em}
\multlinegap0pt
\usepackage{xcolor}
\usepackage{amssymb}
\usepackage{dsfont}
\newtheorem{theorem}{Theorem}
\newcommand{\ip}[2]{\langle #1, \, #2 \rangle}
\newcommand{\mD}{\mbb{D}}
\newcommand{\mN}{\mbb{N}}
\newcommand{\mbb}[1]{\mathds{#1}} % for the consistency of fonts
\newcommand{\mr}[1]{\mathrm{#1}}
\title{Koopman-based Estimation of Lyapunov Functions: \\ Theory on a Reproducing Kernel Hilbert Space}
\author{Wentao Tang, Xiuzhen Ye}
\date{Source: arXiv:2603.01403v1 (CC BY 4.0)}
\begin{document}
\maketitle

\noindent\emph{Source and licence.} Koopman-based Estimation of Lyapunov Functions: \\ Theory on a Reproducing Kernel Hilbert Space. Version arXiv:2603.01403v1 (CC BY 4.0), distributed by arXiv under the Creative Commons Attribution 4.0 International licence, http://creativecommons.org/licenses/by/4.0/, as recorded in the arXiv licence metadata for this exact version and shown on the abstract page https://arxiv.org/abs/2603.01403v1. Full licence notice: \texttt{licenses/arxiv-2603.01403-CC-BY-4.0.txt}.\par\smallskip

\noindent\textit{Editorial scope.} Counted unit: the article's theorem theorem (source0.tex lines 271--273). The article's complete original proof of it is delivered with it (source0.tex lines 274--284, one \texttt{proof} environment of 11 lines). No mathematics is written, repaired or re-derived: the statement and the proof are the article's own lines, in the article's own notation, and no retained line is substituted or re-flowed.  Retained uncounted as the source-local context the proof needs: lines 227--230; lines 258--262; lines 264--266.  Nothing was omitted from any retained span, so the package carries no omission record.  No retained line contains a citation, so the wrapper carries no bibliography and no bibliography entry.  No retained line contains a figure environment, an image-inclusion command or drawing code, so no image is delivered and the package declares no asset dependency.  Editorial formatting only, in addition: an AMS \texttt{amsart} preamble that restates the article's own declarations and macros used by the retained lines, namely the environments \texttt{theorem} on separate counters, together with the standard packages those lines need.  The retained environments print in this document as Theorem 1 in order of appearance, whereas the article prints the same items with the numbering of its own full document.  The complete native source of the article as distributed in the arXiv e-print for this exact version is bundled unchanged as \texttt{sources/195560/source0.tex}.
    \begin{equation}\label{eq:HS}
        \textstyle
        Q = \sum_{i\in\mN} q_iu_i\times u_i
    \end{equation}

\begin{equation}\label{eq:P.series}
    \textstyle 
    v(x) = \ip{\kappa(x, \cdot)}{P\kappa(x, \cdot)}, \enspace 
    P = \sum_{t=0}^\infty A^t Q A^{*t}.
\end{equation}

\begin{equation}\label{eq:P.ale}
    APA^\ast - P = -Q. 
\end{equation}

\begin{theorem}
    Let $H$ be a Hilbert space, $Q\in \mr{HS}(H)$ a positive operator, and $A: H\rightarrow H$ a bounded linear operator with $r(A)<1$, i.e., $\sigma(A)\subset \mD$. Then the the series given in \eqref{eq:P.series} determines a $P\in \mr{HS}(H)$ as a solution to the operator ALE \eqref{eq:P.ale} and in fact the unique solution.  
\end{theorem}

\begin{proof}
    Let $\{u_i\}_{i\in \mN}$ be an orthonormal basis of $H$ and $Q$ be represented as \eqref{eq:HS}. Then for $t\in \mN$, we have
    \begin{align*} 
    \textstyle 
        &\|A^tQA^{\ast t}\|_{\mr{HS}}^2 = \|\sum_{i\in \mN} q_i A^tu_i\times A^tu_i \|_{\mr{HS}}^2 \\
        &= \sum_{i\in \mN}\sum_{j\in \mN} q_iq_j|\ip{A^tu_i}{A^tu_j}|^2 \leq \|Q\|_{\mr{HS}}^2\|A^t\|^4. 
    \end{align*}
    and hence $\|A^tQA^{\ast t}\|_{\mr{HS}} \leq \|Q\|_{\mr{HS}}\|A^t\|^2$. Given any $\epsilon \in (0, (1-r(A))/2)$, there exists a $\overline{t}\in \mN$ such that $\|A^t\|\leq (r(A)+\epsilon)^t$ for all $t\geq \overline{t}$, and hence $\sum_{t=\overline{t}}^\infty A^tQA^{\ast t}$ is an HS operator with HS norm bounded by $\|Q\|_{\mr{HS}} \sum_{t=\overline{t}}^\infty (r(A)+\epsilon)^{2t} \leq (r(A)+\epsilon)^{-2\overline{t}}$. 
    Clearly, $\sum_{t=0}^{\overline{t}-1} A^tQA^{\ast t}$ is HS. Thus the series $P$ is HS. To prove uniqueness, suppose that $P$ and $P'$ are two solutions to \eqref{eq:P.ale}, then $\tilde{P} = P-P'$ satisfies $A\tilde{P}A^\ast = \tilde{P}$. It follows that $\tilde{P} = A^t\tilde{P}A^{\ast t}$ for all $t\in \mN$. Hence we have, 
    $\|\tilde{P}\|_{\mr{HS}} \leq \limsup_{t\rightarrow \infty} \|A^t\tilde{P}A^{\ast t}\|_{\mr{HS}} \leq \|\tilde{P}\|_{\mr{HS}}\limsup_{t\rightarrow \infty} \|A^t\| = 0.$ This implies $\tilde{P}=0$. Therefore the solution to \eqref{eq:P.ale} must be unique. 
\end{proof}

\end{document}
